\documentclass[10pt,a4paper]{amsart}
\usepackage[margin=19mm]{geometry}
\usepackage{amsmath,amssymb,mathtools}
\usepackage[scr=rsfs,cal=euler]{mathalfa}
\usepackage{xfrac}
\usepackage[unicode,hidelinks,bookmarks=false]{hyperref}
\newcommand{\ie}{i.e.\ }
\newcommand{\cf}{cf.\ }
\newcommand{\resp}{respectively\ }
\newcommand{\Subsection}{\subsection}
\providecommand{\nobreakdash}{\nobreak\hbox{-}}
\setlength{\emergencystretch}{2em}
\multlinegap0pt
\usepackage{siunitx}
\usepackage{siunitx}
\usepackage{siunitx}
\usepackage{siunitx}
\usepackage{xcolor}
\usepackage{amssymb}
\usepackage[normalem]{ulem}
\usepackage{stackengine}
\usepackage{siunitx}
\usepackage{tensor}
\usepackage{tikz}
\usepackage[shortlabels]{enumitem}
\usepackage{mathtools}
\newtheorem{definition}{Definition}
\newtheorem{theorem}{Theorem}
\newtheorem{lemma}{Lemma}
\newtheorem{corollary}{Corollary}
\usepackage{cleveref}
\crefname{definition}{Definition}{Definitions}
\crefname{theorem}{Theorem}{Theorems}
\crefname{lemma}{Lemma}{Lemmas}
\crefname{corollary}{Corollary}{Corollarys}
\newcommand{\RR}{{\mathbb{R}}}
\newcommand{\activeS}[2]{A^{#1}_{#2}}
\newcommand{\basis}[1]{\mathcal{#1}}
\newcommand{\bsp}{B}
\newcommand{\bspB}{\basis{\bsp}}
\newcommand{\cU}{\theta}
\newcommand{\cV}{\rho}
\newcommand{\complex}[1]{\mathit{\vSpace{#1}^\bullet}}
\newcommand{\dcU}{\mathit{d\cU}}
\newcommand{\dcV}{\mathit{d\cV}}
\newcommand{\degU}{p^\cU}
\newcommand{\degV}{p^\cV}
\newcommand{\hBasis}[1]{\mathit{\basis{H}\basis{#1}}}
\newcommand{\hSpace}[1]{\mathit{\vSpace{H}\vSpace{#1}}}
\newcommand{\hintComplex}{\complex{H\bsp}^{,\tint}}
\newcommand{\isomorphic}{\cong}
\newcommand{\jet}[2]{J^{#1}\left(#2\right)}
\newcommand{\kntU}{\cU}
\newcommand{\kntV}{\cV}
\newcommand{\kntsU}{\mbf{\kntU}}
\newcommand{\kntsV}{\mbf{\kntV}}
\newcommand{\mbf}[1]{{\boldsymbol{#1}}}
\newcommand{\per}{\mathrm{per}}
\newcommand{\pkntsU}{\kntsU^\per}
\newcommand{\ractiveS}[2]{\tilde{A}^{#1}_{#2}}
\newcommand{\spaceSpan}[1]{\ensuremath{\text{span}}\left(#1\right)}
\newcommand{\tint}{\mathrm{int}}
\newcommand{\tpB}[2]{\bspB^{#1}}
\newcommand{\tpDomain}{\widehat{\Omega}}
\newcommand{\vSpace}[1]{\mathbb{#1}}
\title{Adaptively-refinable polar-spline discrete differential forms:\\ hierarchical construction, exactness, and applications}
\author{Diogo C. Cabanas, Deepesh Toshniwal, Rafael Vazquez}
\date{Source: arXiv:2609.03461v1 (CC BY 4.0)}
\begin{document}
\maketitle

\noindent\emph{Source and licence.} Adaptively-refinable polar-spline discrete differential forms:\\ hierarchical construction, exactness, and applications. Version arXiv:2609.03461v1 (CC BY 4.0), distributed by arXiv under the Creative Commons Attribution 4.0 International licence, http://creativecommons.org/licenses/by/4.0/, as recorded in the arXiv licence metadata for this exact version and shown on the abstract page https://arxiv.org/abs/2609.03461v1. Full licence notice: \texttt{licenses/arxiv-2609.03461-CC-BY-4.0.txt}.\par\smallskip

\noindent\textit{Editorial scope.} Counted unit: the article's lemma lem:intermediate-reduced-complex (source0.tex lines 2335--2337). The article's complete original proof of it is delivered with it (source0.tex lines 2338--2389, one \texttt{proof} environment of 52 lines). No mathematics is written, repaired or re-derived: the statement and the proof are the article's own lines, in the article's own notation, and no retained line is substituted or re-flowed.  Retained uncounted as the source-local context the proof needs: lines 1547--1565; lines 1666--1670; lines 2158--2163; lines 2221--2247; lines 2249--2254.  Nothing was omitted from any retained span, so the package carries no omission record.  No retained line contains a citation, so the wrapper carries no bibliography and no bibliography entry.  No retained line contains a figure environment, an image-inclusion command or drawing code, so no image is delivered and the package declares no asset dependency.  Editorial formatting only, in addition: an AMS \texttt{amsart} preamble that restates the article's own declarations and macros used by the retained lines, namely the environments \texttt{definition}, \texttt{theorem}, \texttt{lemma}, \texttt{corollary} on separate counters, together with the standard packages those lines need.  The retained environments print in this document as Definition 1, Theorem 1, Lemma 1, Corollary 1 in order of appearance, whereas the article prints the same items with the numbering of its own full document.  The complete native source of the article as distributed in the arXiv e-print for this exact version is bundled unchanged as \texttt{sources/209300/source0.tex}.
\begin{definition}[Reduced complex at the pole]\label{def:reduced-complex}
    We say that a complex $\complex{V}$ of forms (with sufficient regularity) on $\tpDomain$ is reduced at the pole if the spaces in it can be decomposed as:
    \begin{equation}
        \begin{split}
            \vSpace{V}^0 &= \RR \oplus g(\cV)\vSpace{V}_\cU^0 \oplus \check{\vSpace{V}}^0\;,\\
            \vSpace{V}^1 &= \left(g(\cV)\vSpace{V}_\cU^1 \oplus \check{\vSpace{V}}^{11}\right)\dcU
                            + \left(\vSpace{V}_\cU^0 1_\cV \oplus \check{\vSpace{V}}^{12}\right)\dcV\;,\\
            \vSpace{V}^2 &= \left(\vSpace{V}_\cU^1 1_\cV \oplus \check{\vSpace{V}}^2\right)\dcU \wedge \dcV\;,
        \end{split}
    \end{equation}
    where:
    \begin{itemize}
        \item $g$ is a $C^1$-smooth function of $\cV$ such that $g(0) = 0$ and $g'(0) \neq 0$;
        \item $\vSpace{V}_\cU^0$ and $\vSpace{V}_\cU^1$ are univariate spaces in $\cU$ such that $\frac{d\vSpace{V}_\cU^0}{\dcU} \subset \vSpace{V}_\cU^1$;
        \item $1_\cV$ is the space of functions constant in $\cV$;
        \item $f \in \check{\vSpace{V}}^0$ or $f \in \check{\vSpace{V}}^{11}$ implies that $\jet{1}{f} = 0$;
        \item $f \in \check{\vSpace{V}}^{12}$ or $f \in \check{\vSpace{V}}^2$ implies that $\jet{0}{f} = 0$.
    \end{itemize}
\end{definition}

\begin{theorem}\label{thm:intermediate-complex}
    The complex $\complex{\bsp}^{,\tint}$ is reduced at the pole.
    Moreover, we have:
    $H^0(\complex{{\bsp}}^{,\tint}) \isomorphic \RR$, $H^1(\complex{{\bsp}}^{,\tint}) = 0$, and $H^2(\complex{{\bsp}}^{,\tint}) = 0$.
\end{theorem}

\begin{lemma}\label{lem:active-intermediate-basis}
    There is a unique level $0 \leq \ell \leq L$ such that $\bsp^{0,\tint}_{1,\ell}$ is active.
    Moreover, except the function identified above, every other active function in $\hBasis{\bsp}^{k,\tint}$ is identical to some level-$\ell'$ tensor-product B-spline $k$-form in $\bspB^{k,\per}_{\ell'}$,
    % $\tpB{k}{\pkntsU_{\ell'},\kntsV_{\ell'}}$,
     $k = 0, 1, 2$.
\end{lemma}

\begin{equation}\label{eq:reduced-active-sets}
    \small
    \begin{split}
        \ractiveS{0,\bsp^\tint}{} &:= \left\{
            (i,\ell') \in \activeS{0,\bsp^\tint}{} ~:~ \jet{1}{\bsp^{0,\tint}_{i,\ell'}} \neq 0
        \right\}\backslash \{ (1,\ell) \}\;,\\
        \ractiveS{1,\bsp^\tint}{} &:= \left\{
            (i,\ell') \in \activeS{1,\bsp^\tint}{} ~:~
            \bsp^{1,\tint}_{i,\ell'} = \bsp^{1,\tint}_{i,\ell',\cU}\;\dcU
            \text{ and }
            \jet{1}{\bsp^{1,\tint}_{i,\ell',\cU}} \neq 0
        \right\} \\
        &\qquad\bigcup
        \left\{
            (i,\ell') \in \activeS{1,\bsp^\tint}{} ~:~
            \bsp^{1,\tint}_{i,\ell'} = \bsp^{1,\tint}_{i,\ell',\cV}\;\dcV
            \text{ and }
            \jet{0}{\bsp^{1,\tint}_{i,\ell',\cV}} \neq 0
        \right\}\;,\\
        \ractiveS{2,\bsp^\tint}{} &:= \left\{
            (i,\ell') \in \activeS{2,\bsp^\tint}{} ~:~
            \bsp^{2,\tint}_{i,\ell'} = \bsp^{2,\tint}_{i,\ell',\cU\cV}\;\dcU\wedge\dcV
            \text{ and }
            \jet{0}{\bsp^{2,\tint}_{i,\ell',\cU\cV}} \neq 0
        \right\}\;.
    \end{split}
\end{equation}

\begin{corollary}\label{cor:reduced-intermediate-basis}
    Given Assumptions 1-2 and the unique level $\ell$ in \cref{lem:active-intermediate-basis}, the sets $\ractiveS{0,\bsp^\tint}{}$, $\ractiveS{1,\bsp^\tint}{}$, and $\ractiveS{2,\bsp^\tint}{}$ are such that:
    \begin{equation*}
        (i, \ell') \in \ractiveS{k,\bsp^\tint}{} \Longrightarrow \ell' \geq \ell\;.
    \end{equation*}
\end{corollary}

\begin{lemma}\label{lem:intermediate-reduced-complex}
    The complex $\hintComplex$ is reduced at the pole.
\end{lemma}

\begin{proof}	
    Let us define the univariate B-spline spaces, $\hSpace{V}_\cU^0$ and $\hSpace{V}_\cU^1$, as follows:
    \begin{equation}\label{eq:reduced-spaces-at-pole}
        \small
        \begin{split}
            \hSpace{V}_\cU^0 &:= \spaceSpan{
                \left\{
                    \bsp_{i_1} : \bsp^{0,\tint}_{i,\ell'} = \bsp_{i_1,\degU,\ell}\bsp_{i_2,\degV,\ell}\text{ for some }(i,\ell') \in \ractiveS{0,\bsp^\tint}{}\;
                \right\}
            }\;,\\
            \hSpace{V}_\cU^1 &:= \spaceSpan{
                \left\{
                    \bsp_{i_1} : \bsp^{1,\tint}_{i,\ell'} = \bsp_{i_1,\degU,\ell}\bsp_{i_2,\degV,\ell}\;\dcU\text{ for some }(i,\ell') \in \ractiveS{1,\bsp^\tint}{}
                \right\}
            }\;,
        \end{split}
    \end{equation}
    where we recall that $\ractiveS{0,\bsp^\tint}{}$ and $\ractiveS{1,\bsp^\tint}{}$ are defined in \cref{eq:reduced-active-sets}.
    Moreover, and analogously to \cref{thm:intermediate-complex}, we define the spaces:
    \begin{equation*}
        \begin{split}
            \check{\vSpace{V}}^0 &:= \spaceSpan{B~:~B \in \hBasis{\bsp}^{0,\tint}
            ~~\text{and}~~
            \jet{1}{B} = 0}\;,\\
            \check{\vSpace{V}}^{11} &:= \spaceSpan{B~:~B\;\dcU \in \hBasis{\bsp}^{1,\tint}
            ~~\text{and}~~
            \jet{1}{B} = 0}\;,\\
            \check{\vSpace{V}}^{12} &:= \spaceSpan{B~:~B\;\dcV \in \hBasis{\bsp}^{1,\tint}
            ~~\text{and}~~
            \jet{0}{B} = 0}\;,\\
            \check{\vSpace{V}}^2 &:= \spaceSpan{B~:~B\;\dcU\wedge\dcV \in \hBasis{\bsp}^{2,\tint}
            ~~\text{and}~~
            \jet{0}{B} = 0}\;.
        \end{split}
    \end{equation*}
    Note that these spaces are different from those in \cref{thm:intermediate-complex}; we reuse the notation simply to establish a connection with \cref{def:reduced-complex}.
    Then, with $g(\cV) := \bsp_{2,\degV,\ell}(\cV)$ and level $\ell$ as in \cref{lem:active-intermediate-basis}, we have that the spaces $\hSpace{\bsp}^{0,\tint}$, $\hSpace{\bsp}^{1,\tint}$, and $\hSpace{\bsp}^{2,\tint}$ can be expressed in the following form:
    \begin{equation*}
        \begin{split}
            \hSpace{\bsp}^{0,\tint} &= \RR \oplus g(\cV)\hSpace{V}_\cU^0 \oplus \check{\vSpace{V}}^0\;,\\
            \hSpace{\bsp}^{1,\tint} &= \left(g(\cV)\hSpace{V}_\cU^1 \oplus \check{\vSpace{V}}^{11}\right)\dcU
                            + \left(\hSpace{V}_\cU^01_\cV \oplus \check{\vSpace{V}}^{12}\right)\dcV\;,\\
            \hSpace{\bsp}^{2,\tint} &= \left(\hSpace{V}_\cU^1 1_\cV \oplus \check{\vSpace{V}}^2\right)\dcU \wedge \dcV\;.
        \end{split}
    \end{equation*}
    
    For example, consider $\hSpace{\bsp}^{0,\tint}$ and let $B \in \hBasis{\bsp}^{0,\tint}$ such that $\jet{0}{B} = 0$ and $\jet{1}{B} \neq 0$.
    With \cref{cor:reduced-intermediate-basis} in mind, $B$ must be a B-spline from level $\ell' \geq \ell$.
    By nestedness of the level-wise spaces (particularly in the coordinate $\cV$), $B$ can be expressed as $B_1 + B_2$ where $B_1 \in g(\cV)\hSpace{V}_\cU^0$ with $\jet{1}{B_1} = \jet{1}{B}$, and $B_2 \in \check{\vSpace{V}}^0$ with $\jet{1}{B_2} = 0$.

    The reasoning is similar for the other spaces and, finally, it is easy to check that \cref{def:reduced-complex} applies.
\end{proof}

\end{document}
